\section{Experimental Setup}

We design experiments to answer four research questions that directly test our representational alignment hypothesis and scaffolding framework:

\begin{itemize}
    \item \textbf{RQ1}: Does structured error formatting improve repair success over raw compiler output?
    \item \textbf{RQ2}: Is the improvement due to organizational structure, or merely information content?
    \item \textbf{RQ3}: What is the optimal level of fix specificity?
    \item \textbf{RQ4}: Does model scale interact with format benefit?
\end{itemize}

\subsection{Datasets}

We evaluate on EvalPlus benchmarks \citep{guo2024evalplus}, which provide rigorous evaluation with 80$\times$ more test cases than original HumanEval/MBPP:

\begin{table}[H]
\centering
\begin{tabular}{@{}llll@{}}
\toprule
\textbf{Dataset} & \textbf{Problems} & \textbf{Tests per Problem} & \textbf{Source} \\
\midrule
HumanEval+ & 164 & $\sim$80 avg & EvalPlus \\
MBPP+ & 378 & $\sim$80 avg & EvalPlus \\
\textbf{Total} & 542 & $\sim$43,360 & -- \\
\bottomrule
\end{tabular}
\caption{Evaluation datasets.}
\end{table}

\textbf{Why EvalPlus}: Original HumanEval/MBPP have saturated (99.4\% and 94.2\% pass@1 for top models). EvalPlus's expanded test suites provide discriminative evaluation and generate more diverse error instances for our experiments.

\subsubsection{Error Collection Process}

\begin{enumerate}
    \item Run base model generation on all problems
    \item Collect failed test cases with static analysis errors
    \item Parse errors to extract: line number, error type, error message, code context
    \item Filter for 8 supported error types (balanced sampling)
\end{enumerate}

We collect 500 error instances per experiment, stratified across error types to ensure balanced representation.

\subsection{Models}

We evaluate across three model scales to test scale interaction:

\begin{table}[H]
\centering
\begin{tabular}{@{}llll@{}}
\toprule
\textbf{Model} & \textbf{Parameters} & \textbf{Access} & \textbf{Purpose} \\
\midrule
CodeLlama-7B-Instruct & 7B & HuggingFace & Small scale \\
CodeLlama-34B-Instruct & 34B & HuggingFace & Medium scale \\
GPT-4 & $\sim$175B+ & OpenAI API & Large scale \\
\bottomrule
\end{tabular}
\caption{Models evaluated.}
\end{table}

\textbf{Why this selection}: CodeLlama variants provide controlled comparison within a model family, while GPT-4 represents current state-of-the-art for code generation. This span enables testing our hypothesis that smaller models benefit more from structured formatting.

\subsection{Baselines}

\textbf{Primary baseline}: Self-repair with raw compiler output, following the theoxo/self-repair framework \citep{theoxo2024iclr}. This represents the standard practice of feeding raw tracebacks directly to the model.

\textbf{Control conditions}:
\begin{itemize}
    \item \textbf{Verbose-Raw}: Expanded raw output with additional whitespace and formatting---controls for prompt length
    \item \textbf{Scrambled}: Structured format content with randomly permuted section order---controls for information content
\end{itemize}

\textbf{Why these baselines}: The Scrambled condition is critical for causal inference. If Structured $>$ Scrambled, the effect is due to organizational structure, not information availability. Verbose-Raw rules out simple length effects.

\subsection{Sub-Hypothesis Experiments}

\subsubsection{H-E1: Existence Test (MUST\_WORK)}
\begin{itemize}
    \item \textbf{Design}: Within-subject comparison of Raw vs Structured format
    \item \textbf{N}: 500 error instances $\times$ 2 conditions
    \item \textbf{Metric}: Repair success rate (binary)
    \item \textbf{Success criterion}: Structured $>$ Raw with $p < 0.05$
\end{itemize}

\subsubsection{H-M1: Information Preservation (MUST\_WORK)}
\begin{itemize}
    \item \textbf{Design}: Reconstruction test using third-party LLM
    \item \textbf{N}: 500 error pairs across 8 error types
    \item \textbf{Metric}: Field extraction accuracy
    \item \textbf{Success criterion}: $>$95\% reconstruction accuracy
\end{itemize}

\subsubsection{H-M2: Representational Alignment (SHOULD\_WORK)}
\begin{itemize}
    \item \textbf{Design}: Paired comparison of Structured vs Scrambled
    \item \textbf{N}: 500 error instances, paired
    \item \textbf{Statistical test}: McNemar's chi-squared for paired binary outcomes
    \item \textbf{Success criterion}: Structured $>$ Scrambled with $p < 0.05$
\end{itemize}

\subsubsection{H-M3: Fix Specificity (SHOULD\_WORK)}
\begin{itemize}
    \item \textbf{Design}: 4-level within-subject comparison (Levels 0-3)
    \item \textbf{N}: 500 instances $\times$ 4 levels $\times$ 3 repetitions
    \item \textbf{Statistical test}: Mixed-effects model with quadratic contrast
    \item \textbf{Success criterion}: Significant quadratic term, peak at Level 1-2
\end{itemize}

\subsubsection{H-C1: Scale Interaction (SHOULD\_WORK)}
\begin{itemize}
    \item \textbf{Design}: 2$\times$3 factorial (Format $\times$ Model Scale)
    \item \textbf{N}: 542 problems $\times$ 6 cells
    \item \textbf{Statistical test}: Two-way ANOVA with interaction term
    \item \textbf{Success criterion}: Interaction $p < 0.05$, $\eta^2 > 0.01$
\end{itemize}

\subsection{Evaluation Metrics}

\textbf{Primary metric}: Repair success rate---proportion of errors successfully fixed within max iterations.

\textbf{Secondary metrics}:
\begin{itemize}
    \item Pass@1 improvement: Delta vs base model without repair
    \item Iterations to success: Attempts required (when successful)
    \item Error type breakdown: Success rate per error category
\end{itemize}

\textbf{Statistical significance}: All comparisons use $p < 0.05$ threshold with Benjamini-Hochberg FDR correction for multiple comparisons. Effect sizes reported as Cohen's d for continuous outcomes and odds ratios for binary outcomes.

\subsection{Implementation Details}

\textbf{Inference}: Temperature = 0.0 for deterministic comparison. Max repair iterations = 5. Batch size = 1 (sequential repair).

\textbf{Hardware}: 5$\times$ NVIDIA H100 NVL (95GB each). CodeLlama models run locally; GPT-4 via API.

\textbf{Reproducibility}: Seeded randomization (seed=42) for all stochastic components. All code available in supplementary materials.
